\documentclass[../relazione.tex]{subfiles}
\begin{document}

% ======================================================================
\section{Discussione}
% ======================================================================

\paragraph{Implicazioni principali:}
L'analisi comparativa su 5 modelli Pythia (160M -- 6.9B) permette di trarre conclusioni significative:

\begin{enumerate}
    \item \textbf{Il bias liberal è intrinseco ai dati ma è transitorio}: Pythia-160M mostra D+L = 50.4\%, ma il bias si riduce con la scala. A partire da 1.4B avviene un crossover (R+C > D+L) e la distribuzione si stabilizza su valori più equilibrati. Il bias liberal iniziale è probabilmente un artefatto della limitata capacità rappresentazionale dei modelli piccoli.

    \item \textbf{Lo scaling migliora la robustezza con rendimenti decrescenti}: il salto qualitativo avviene tra 160M e 1B (Position Bias: 78\% $\to$ 5\%). Da 1.4B in poi, i miglioramenti sono marginali ($>94\%$ Robust per tutti).

    \item \textbf{L'allineamento umano ha un plateau}: l'alignment cresce fino a 1.4B (0.821) poi si stabilizza ($\sim$0.81--0.82). Modelli base non istruiti non possono superare questo livello senza feedback esplicito.

    \item \textbf{Le minacce hanno effetti non monotoni}: il validity rate sotto threat non cresce linearmente con la scala. Il picco si raggiunge a 2.8B (54.1\%), con un calo a 6.9B (44.9\%). Contemporaneamente, il tasso di Collapses cresce (8.2\% a 6.9B).

    \item \textbf{La minaccia più destabilizzante evolve}: da Economic (modelli piccoli) a Legal (modelli grandi), riflettendo una crescente capacità di ``comprendere'' la gravità delle conseguenze.

    \item \textbf{IT/System è la più efficace per formato}: raggiunge validity rate del 72.2\% (2.8B) e 70.2\% (6.9B) e la coerenza logit-testo più alta (0.573 a 6.9B). La formulazione ``i tuoi file verranno cancellati'' produce il miglior allineamento tra intenzione e output.

    \item \textbf{Confronto con Santurkar et al.}: GPT-3 Davinci allineato mostra 89\% Democrat, 7\% Republican. I nostri modelli base mostrano bias molto più moderati, con una convergenza verso l'equilibrio già da 1.4B. Questo conferma che l'RLHF \emph{amplifica drasticamente} il bias preesistente nei dati.
\end{enumerate}

\paragraph{Limitazioni dello studio:}
\begin{itemize}
    \item \textbf{Basso validity rate}: anche 6.9B raggiunge solo il 14.2\% baseline. Le metriche basate sui logit (JSD, WD, alignment) coprono comunque tutte le domande.
    \item \textbf{Solo modelli Pythia}: per generalizzare servirebbero modelli con architetture e corpora diversi.
    \item \textbf{Non-monotonia del threat}: il comportamento a 1.4B e 6.9B suggerisce fattori non catturati dalla sola dimensione (es. architettura, numero di layers).
\end{itemize}

\end{document}
